\subsection{Extremal generators of convex cones}

Let C be a convex cone with vertex 0 in a vector space E; clearly no other point of C than the vertex can be an extremal point; the vertex is an extremal point of C if and only if C is pointed and proper.

DEFINITION 2. --- Let C be a convex cone of vertex 0 in a vector space E. We say that a half-line D ⊂ C originating at 0 is an extremal generator of C, if every open segment contained in C, not containing 0 and meeting D is contained in D.

It comes to the same thing to say that for all \(x \in D\) such that \(x \neq 0\) , if \(y \neq 0\) , \(y' \neq 0\) are two points of C such that \(x = y + y'\) , then, it is necessarily the case that \(y \in D\) and \(y' \in D\) .

Remark 1. --- Let C be a pointed proper convex cone in E, and consider on E the order structure for which C is the set of elements \(\geqslant 0\) (II, p.~12, prop. 13); in order that an element of E, say x \textgreater{} 0, belongs to an extremal generator of C, it is necessary and sufficient that every element \(y \geqslant 0\) , that is bounded above by x, is of the form \(\lambda x\) with \(0 \leqslant \lambda \leqslant 1\) : in fact, to say that y is bounded above by x means that \(x = y + y'\) where \(y' \in C\) , whence the conclusion follows.

PROPOSITION 3. --- In a vector space E, let C be a convex cone with vertex 0, and let \(x_{0} \neq 0\) be a point of C, and D a half-line that is contained in C, originating from 0 and containing \(x_{0}\) . Let H be a hyperplane containing \(x_{0}\) and not passing through 0. Then D is an extremal generator of C if and only if \(x_{0}\) is an extremal point of H ∩ C.

The condition is clearly necessary. Conversely, suppose that it is satisfied; suppose that there is a line \(D'\) not containing D, passing through \(x_{0}\) and such that \(D' \cap C\) contains an open segment to which \(x_{0}\) belongs. Let \(y \neq 0\) be a direction vector of \(D'\) ; the hypotheses imply that the point \((1 + \lambda)x_{0} + \mu y\) belongs to C for \(|\lambda|\) and \(|\mu|\) sufficiently small. But then, in the plane P determined by D and \(D'\) and carrying the canonical topology, \(x_{0}\) is an interior point of \(P \cap C\) , and it follows that the line \(P \cap H\) contains an open segment contained in \(H \cap C\) and to which \(x_{0}\) belongs. This contradicts the hypothesis.

DEFINITION 3. --- Let C be a convex set in a Hausdorff topological vector space E. A compact convex non-empty set A of C is called a cap of C if the complement C−A of A in C is convex.

Let C be a pointed convex cone with vertex 0 in E and let A be a cap of C. Write B = C - A. For every closed half-line L ⊂ C originating at 0, the sets L ∩ A and L ∩ B are convex sets that are complements in L, whose union is L, and such that L ∩ A is compact. As L ∩ A is non-empty for at least one half-line L, we see that 0 ∈ A, thus L ∩ A is a closed segment with an end point at 0. If A exists then C is proper.

PROPOSITION 4. --- Let C be a pointed convex cone with vertex 0 in E, a Hausdorff locally convex space.

\begin{enumerate}
\def\labelenumi{\alph{enumi})}
\tightlist
\item
  Let A be a cap of C. Let p be the restriction to C of the gauge of A (II, p.~20). The set of the \(x \in C\) such that \(p(x) \leqslant 1\) is the set A. The function p is lower semicontinuous and has the following properties :
\end{enumerate}

\begin{enumerate}
\def\labelenumi{(\roman{enumi})}
\item
  For any \(x, y\) in \(\mathbf{C}\) , we have \(p(x + y) = p(x) + p(y)\) .
\item
  For any \(x \in \mathbf{C}\) and \(\lambda \in \mathbf{R}_{+}^{*}\) , we have \(p(\lambda x) = \lambda p(x)\) .
\item
  If \(x \in C\) , the relation \(p(x) = 0\) is equivalent to x = 0.
\end{enumerate}

\begin{enumerate}
\def\labelenumi{\alph{enumi})}
\setcounter{enumi}{1}
\tightlist
\item
  Conversely, let \(p\) be a function defined in \(\mathbf{C}\) with values in \([0, +\infty]\) , satisfying the conditions (i), (ii) of a). Let \(\mathbf{A}\) be the set of the \(x \in \mathbf{C}\) such that \(p(x) \leqslant 1\) . Then \(\mathbf{A}\) and \(\mathbf{C} - \mathbf{A}\) are convex. A is a cap, if and only if \(\mathbf{A}\) is compact and non-empty.
\end{enumerate}

The statement \(b)\) is obvious. The properties stated in \(a)\) are consequences of the remarks preceding prop. 4 and of the prop. 22 of II, p.~20 and prop. 23 of II, p.~20 with the exception of

\[
p (x + y) \geqslant p (x) + p (y).
\]

It is sufficient to prove this last when \(x \neq 0\) and \(y \neq 0\) ; we have therefore \(p(x) > 0\) , \(p(y) > 0\) . Let \(\mu, \lambda\) be two numbers \textgreater{} 0 such that \(\lambda < p(x)\) , \(\mu < p(y)\) , and denote the complement of A in C by B. We have \(x \in \lambda B\) , \(y \in \mu B\) , therefore \(x + y \in \lambda B + \mu B\) ; by the convexity of B, we have \(\lambda B + \mu B \subset (\lambda + \mu)B\) , whence \(p(x + y) > \lambda + \mu\) , which implies the inequality stated above.

COROLLARY 1. --- Let C be a pointed convex cone of vertex 0 in E, a Hausdorff locally convex space and let p be the gauge of A, a cap of C. The extremal points of A are then the point 0, and the points x on the extremal generators of C such that \(p(x) = 1\) .

It is clear that 0 is an extremal point of A. Let x be a point on L an extremal generator of C and such that \(p(x) = 1\) . Let y, z be two points of A such that \(x = \frac{1}{2}(y + z)\) . As L is extremal, we have \(y = \lambda x\) and \(z = \mu x\) , where \(\lambda\) and \(\mu\) are numbers \(\geqslant 0\) such that \(\frac{1}{2}(\lambda + \mu) = 1\) , \(\lambda = \lambda p(x) = p(y) \leqslant 1\) and \(\mu = \mu p(x) = p(z) \leqslant 1\) , from which \(\lambda = \mu = 1\) and hence y = z = x; so, x is an extremal point of A. Conversely, let \(x \neq 0\) be an extremal point of A. Obviously \(p(x) = 1\) . Let y, \(y'\) be two points of C such that \(x = y + y'\) , and we shall show that y, \(y'\) are proportional to x. Without loss of generality we can suppose that the numbers \(\lambda = p(y)\) and \(\lambda' = p(y')\) are finite and \textgreater0. Then \(\lambda^{-1}y \in A\) , \(\lambda'^{-1}y' \in A\) , \(\lambda + \lambda' = 1\) by prop. 4, (i) and the equality \(x = \lambda(\lambda^{-1}y) + \lambda'(\lambda'^{-1}y')\) implies, by hypothesis that

\[
x = \lambda^ {- 1} y = \lambda^ {\prime - 1} y ^ {\prime}.
\]

COROLLARY 2. --- Every point of C that belongs to a cap of C, also belongs to the convex closed envelope of the union of the extremal generators of C.

This follows immediately from cor. 1 and the Krein-Milman theorem (II, p.~55, th. 1).

* Example. --- Let X be a locally compact space that is \(\sigma\) -compact. Let C be a closed convex cone of vertex 0 in \(\mathcal{M}_{+}(X)\) with the vague topology. We shall show that C is the union of its caps. Let \((X_{n})\) be an increasing sequence of open, relatively compact sets of X whose union is X. Let \(\mu\) be an element \(\neq 0\) of C. There exist \(\alpha_{n} > 0\) such that \(\sum \alpha_{n}\mu(X_{n}) = 1\) .

For every measure \(v \in C\) , put \(p(v) = \sum_{n} \alpha_{n} v(X_{n}) \in [0, +\infty]\) . The function p on C satisfies conditions (i) and (ii) of prop. 4. It is lower semi-continuous for the vague topology (INT, IV, 2nd ed., § 1, No.~1, prop. 4). The set A of the \(\gamma \in C\) such that \(p(\gamma) \leqslant 1\) is therefore closed and non-empty. On the other hand, every compact set of X is contained in one of the \(X_{n}\) , thus A being vaguely bounded is also vaguely compact (INT, III, 2nd ed., § 1, No.~9, prop. 15). The set A is therefore a cap of C containing \(\mu\) .

PROPOSITION 5. --- Let C be a proper convex cone with vertex 0 in E, a Hausdorff weak space; suppose that C is complete for the uniform structure induced by that of E, and that there is an enumerable fundamental system of neighbourhoods of 0 in C. Then C is the union of its caps and is the closed convex envelope of the union of its extremal generators.

The second statement follows from the first and from cor. 2 above. Using prop. 11 of II, p.~52 reduces the proposition to the case when \(\mathbf{E} = \mathbf{R}^{\mathrm{I}}\) and \(\mathbf{C} \subset \mathbf{R}_{+}^{\mathrm{I}}\) . For all \(\alpha \in \mathbf{I}\) , denote the projection \(\operatorname{pr}_{\alpha}\) in \(\mathbf{E}\) by \(f_{\alpha}\) ; then \(f_{\alpha}\) is a continuous linear form. On the other hand let \((\mathrm{V}_n)_{n \in \mathbb{N}}\) be an enumerable fundamental system of neighbourhoods of 0 in \(\mathbb{C}\) . By the definition of the topology of \(\mathbf{E}\) , for each \(n \in \mathbb{N}\) , there exists a finite subset \(\mathrm{J}_n\) of \(\mathbf{I}\) and a number \(\varepsilon_n > 0\) such that \(\mathrm{V}_n\) contains the set \(\mathrm{W}_n\) of the \(x \in \mathbb{C}\) such that \(f_{\alpha}(x) \leqslant \varepsilon_n\) for all \(\alpha \in \mathrm{J}_n\) ; put \(\mathrm{J} = \bigcup_{n \in \mathbb{N}} \mathrm{J}_n\) . Let \(y \neq 0\) be a point of \(\mathbb{C}\) , and \(p\) be the function \(\sum_{\alpha \in \mathbf{J}} \lambda_{\alpha}(f_{\alpha}|\mathbb{C})\) where the \(\lambda_{\alpha} > 0\) are chosen so that \(p(y) = 1\) ; this is possible, since if \(f_{\alpha}(y) = 0\) for all \(\alpha \in \mathbf{J}\) , then \(y \in \mathrm{V}_n\) for all \(n\) , which implies that \(y = 0\) , and this is contrary to hypothesis. Now we remark that for all \(\alpha \in \mathbf{I}\) , the function \(f_{\alpha}|\mathbb{C}\) is continuous at the point 0, therefore there is an \(n \in \mathbb{N}\) , such that \(f_{\alpha}\) is bounded in a \(\mathrm{W}_n\) , therefore bounded above in \(\mathbb{C}\) by a linear combination of a finite number of functions \(f_{\beta}|\mathbb{C}\) , where \(\beta \in \mathbf{J}\) . It follows that if \(\mathbf{A}\) in the set of \(x \in \mathbb{C}\) such that \(p(x) \leqslant 1\) , then \(f_{\alpha}\) is bounded in \(\mathbf{A}\) for all \(\alpha \in \mathbf{I}\) . As \(p\) is lower semi-continuous in \(\mathbb{C}\) , it follows that \(\mathbf{A}\) is closed and non-empty in \(\mathbb{C}\) and therefore is compact. Since it is clear that \(p\) verifies the conditions (i) and (ii) of prop. 4 of II, p.~58, we see that \(\mathbf{A}\) is a cap in \(\mathbb{C}\) and contains \(y\) .

Remark 2. --- There exist proper convex cones that are weakly complete and which have no extremal generator (II, p.~92, exerc. 31).

